\section{What the count means for a single equation}\label{sec:fixed}

The count lets the operator vary. We now read it for one evolution
equation \eqref{eq:pde} with fixed complex coefficients. The equation is
invariant under translations and under Galilean shifts.

\begin{proposition}[The waves of a fixed equation]\label{prop:fixed}
Let \(\kappa\ne0\). Up to a translation and a Galilean shift, the
travelling waves of \eqref{eq:pde} that are rational in
\(e^{\kappa(x-ct)}\), with one pole modulo
\((2\pi i/\kappa)\ZZ\), are
\begin{equation}\label{eq:fixed-wave}
 u=b_p\kappa^p\,v_A\bigl(\kappa(x-ct)\bigr),\qquad
 c=-b_p\kappa^pP_A(0),
\end{equation}
where \(A\) ranges over the points of \(\Xp\) that satisfy
\begin{equation}\label{eq:fixed-conditions}
 \frac{b_j}{b_p}=\kappa^{p-j}\,[\rho^j]P_A,\qquad 1\le j\le p-1 .
\end{equation}
The data \((A,\kappa)\) and \((\iota A,-\kappa)\) describe the same
wave, read in the two directions (Section~\ref{sec:reflection}).
\end{proposition}

\begin{proof}
Such a wave has a finite limit as \(e^{\kappa(x-ct)}\to0\)
(Section~\ref{sec:systems}), and a Galilean shift makes this limit
zero. Then \(w_\ast=0\) in \eqref{eq:scaling}, and
\eqref{eq:coefficients} becomes \eqref{eq:fixed-conditions} together
with the formula for \(c\). After a translation \(v\) has the
normalization of Section~\ref{sec:setting}, so \(v=v_A\) and \(P=P_A\)
by Proposition~\ref{prop:divisibility}. The last assertion follows from
\eqref{eq:reflection-pair}: the two waves differ by the constant
\(b_p\kappa^ps_pA(0)\), which is a Galilean shift.
\end{proof}

Given a pair, \eqref{eq:fixed-conditions} are \(p-1\) conditions on the
\(p-1\) coefficient ratios \(b_j/b_p\) and the one unknown \(\kappa\):
a pair imposes \(p-2\) conditions on the equation and fixes the rate.
A pair with real coefficients gives a real wave of a real equation,
smooth on the real axis, exactly when \(\kappa\) is real.
Let \(\mathbf a(A)=([\rho^{p-1}]P_A,\ldots,[\rho^{1}]P_A)\) be the
vector of lower coefficients of \(P_A\), and let
\(\mathbf b=(b_{p-1}/b_p,\ldots,b_1/b_p)\) be the coefficient vector of
the equation. Condition \eqref{eq:fixed-conditions} says that
\(\mathbf b\) is obtained from \(\mathbf a(A)\) by multiplying its
\(k\)th entry by \(\kappa^k\), for some \(\kappa\ne0\). As \(\kappa\)
varies these vectors trace a curve, denoted
\(\CC^\ast\mathbf a(A)\), and rescaling \(x\) moves an equation along
it. This is the reason for fixing the rate: without it every pair would
be a one-parameter family. Each of the \(N_p\) pairs marks one curve,
and several pairs may mark the same curve.

\begin{corollary}[Solvable equations]\label{cor:solvable}
Assume that the matching equations have finitely many solutions.
This holds when \(2p-1\) is prime, and whenever rigidity holds
(Proposition~\ref{prop:bezout}), hence for every \(p\le573\)
(Section~\ref{sec:further}).
\begin{enumerate}
\item[\textup{(i)}] Equation \eqref{eq:pde} has a travelling
wave that is rational in an exponential, with one pole per period, if
and only if \(\mathbf b\) lies on one of the curves
\(\CC^\ast\mathbf a(A)\), \(A\in\Xp\). There are at most
\((N_p+F_p)/2\) such curves.
\item[\textup{(ii)}] If \(\mathbf b\ne0\), there are finitely many such
waves up to translations and Galilean shifts.
\item[\textup{(iii)}] If \(p\ge3\), the curves form a set of dimension
at most one in \(\CC^{p-1}\). A generic equation
\eqref{eq:pde} has no such wave.
\end{enumerate}
\end{corollary}

\begin{proof}
Part~(i) restates Proposition~\ref{prop:fixed}. For the bound,
\(\mathbf a(\iota A)=(-1)\cdot\mathbf a(A)\) by
\eqref{eq:reflection-pair}, so \(A\) and its mirror image \(\iota A\)
give the same curve. The number of classes \(\{A,\iota A\}\) is half
the number of points plus half the number of fixed points. The
solutions are finite in number, hence isolated, so these two numbers
are at most \(N_p\) and \(F_p\) (Proposition~\ref{prop:bezout},
Remark~\ref{rem:Fp-bound}). For~(ii), if \(\mathbf b\ne0\) and
\(\mathbf b\in\CC^\ast\mathbf a(A)\), some coordinate \([\rho^j]P_A\)
is nonzero, and \eqref{eq:fixed-conditions} leaves at most \(p-j\)
values of \(\kappa\). Part~(iii) is a dimension count.
\end{proof}

When \(\mathbf b=0\) the equation is invariant under the scaling itself,
and its waves come in one-parameter families of rates. At order two
there are two curves, and both carry waves: \(\mathbf b=0\) is the KdV
equation, with its solitons of all rates, and \(\mathbf b\ne0\) consists
of the equations with a nonzero Burgers coefficient, each with one
family of fronts, at the rate fixed by that coefficient
(Section~\ref{sec:kdvb}).

At \(p=3\), write \(P=D^3+BD^2+\mu D+a_0\) and use \((B,\mu)\) as the
coordinates on \(\CC^2\). Rescaling multiplies \(B\) by \(\kappa\) and
\(\mu\) by \(\kappa^2\), so the curves are \(B^2=\sigma^2\mu\) with the
origin removed.
For the ten pairs \(A=\rho^2+a\rho+b\) of Table~\ref{tab:ks},
\eqref{eq:ksoperator} gives \(B=4a\) and \(\mu=a^2+19b\), hence the
four classical values
\begin{equation}\label{eq:ks-sigma}
 \sigma^2=\frac{16a^2}{a^2+19b}\in
 \{0,\ 16,\ 144/47,\ 256/73\},
 \qquad\text{that is,}\qquad
 |\sigma|\in\{0,\ 4,\ 12/\sqrt{47},\ 16/\sqrt{73}\},
\end{equation}
which are those of \citet[Eq.~(39)]{Kudryashov1989}. In the
normalization of the Introduction, with \(b_1=b_3=1\) and
\(b_2=\sigma\), condition \eqref{eq:fixed-conditions} reads
\(\kappa^2=1/(a^2+19b)\) and \(\sigma=4a\kappa\).
Here the bound \((N_3+F_3)/2=6\) is not attained, because the two
pairs with \(a=0\) share \(\sigma=0\), and \((\pm1,0)\) and
\((\pm i,0)\) share \(\sigma^2=16\).
For coefficients off these curves, approximate solutions close to the
exact ones are constructed by \citet{KudryashovKochanov2012}.

For odd \(p\) the classification turns
Corollary~\ref{cor:solvable} into a statement about all meromorphic
travelling waves.

\begin{corollary}[Generic equations with \(p\) odd]
\label{cor:generic}
Let \(p\ge5\) be odd and assume rigidity \((\mathrm R_{p-1})\), which
holds when \(2p-1\) is prime and for every \(p\le573\). The coefficient
vectors \(\mathbf b\in\CC^{p-1}\) for which \eqref{eq:pde} has a
nonconstant meromorphic travelling wave lie in an algebraic subset of
dimension at most two. In particular a generic equation has none.
\end{corollary}

\begin{proof}
By Theorem~\ref{thm:complete}, such a wave is rational, rational in an
exponential, or elliptic, with one pole per period. Take
\(\kappa=1\) in \eqref{eq:scaling}, so that
\([\rho^j]P=b_j/b_p\) for \(1\le j\le p-1\).
The rational-exponential waves give a set of dimension at most one, by
Corollary~\ref{cor:solvable}. A rational wave is, after a translation,
of the form \eqref{eq:rational-system}. By Lemma~\ref{lem:laplace} it
comes from a monic \(A\) of degree \(p-1\) that divides \(H_A\). By
rigidity, \(A=\rho^{p-1}\), hence \(P=D^p\) and \(\mathbf b=0\).
For the elliptic waves, Steps~1--3 of the proof of
Theorem~\ref{thm:elliptic-count} use only rigidity, and with them
Lemma~\ref{lem:conservation}(i) shows that every system
\(\mathcal E_p(g_2,g_3)\) has finitely many solutions. The solutions
for all \((g_2,g_3)\) together therefore form an algebraic set of
dimension at most two, and so does the closure of its image under
\(P\mapsto\mathbf b\).
\end{proof}

At \(p=3\) the argument gives the exact answer: the set of these
\(\mathbf b=(B,\mu)\) is the union of the four curves
\(B^2=\sigma^2\mu\) of \eqref{eq:ks-sigma}. Indeed the elliptic waves
require \(B^2=16\mu\) (Section~\ref{sec:ks}), which is one of the four
curves, and the point \(\mathbf b=0\) of the rational wave lies on all
of them.

The same dimension count applies to the one-pole waves of every order
\(p\ge3\) for which rigidity holds: a generic equation has no
travelling wave that is rational, rational in an exponential, or
elliptic with one pole per period. For even \(p\), however,
Theorem~\ref{thm:complete} excludes other meromorphic waves only when
the equation is purely dispersive.

Can one equation have two different waves of this kind at the same
rate and speed? It cannot, and this holds for every \(P\), including
the mixed equations of odd order, to which Theorem~\ref{thm:complete}
does not apply: the argument only uses the leading coefficient of a
convolution.

\begin{proposition}[One profile per operator]\label{prop:injective}
Let \(A_1,A_2\) be monic polynomials of degree \(p-1\) satisfying
\(A_1\mid S_{A_1}\) and \(A_2\mid S_{A_2}\). If \(P_{A_1}=P_{A_2}\),
then \(A_1=A_2\).
Consequently a monic operator of degree \(p\) admits at most one
normalized rational-exponential one-pole profile.
Any count of distinct pairs with this normalization is therefore
also a count of distinct operators.
\end{proposition}

\begin{proof}
Suppose \(P_{A_1}=P_{A_2}=P\), and put \(h=A_1-A_2\).
If \(h\ne0\), its degree \(n\) is less than \(p-1\); write
\(f\ne0\) for its leading coefficient.
By \eqref{eq:operator-A} and the bilinearity and symmetry of the
convolution, \(S_{A_1}-S_{A_2}=h\star(A_1+A_2)=(2/s_p)Ph\).
The power-sum calculation in \eqref{eq:Sleading} gives the leading
coefficient \(2f\int_0^1t^n(1-t)^{p-1}\dd t\) on the left.
Since \(P\) is monic, the right side has leading coefficient
\((2/s_p)f\), and \(2/s_p\) is the beta integral of
\eqref{eq:Sleading}. Cancelling \(f\) would therefore give
\[
 2\int_0^1t^n(1-t)^{p-1}\dd t
   =\int_0^1t^{p-1}(1-t)^{p-1}\dd t.
\]
This is impossible: \(n<p-1\) implies \(t^n>t^{p-1}\) for \(0<t<1\), so
the left side exceeds twice the right side. Thus \(h=0\).
\end{proof}

Every \(P_A\) has a root among \(1,\ldots,p\). Let \(m\) be the least
positive integer with \(A(m)\ne0\); then \(m\le p\), because \(A\) has
degree \(p-1\). As \(e^z\to0\) the profile begins with
\(s_p(-1)^{m-1}A(m)e^{mz}\) and its square with \(e^{2mz}\), so the
coefficient of \(e^{mz}\) in \eqref{eq:main} gives \(P_A(m)=0\).
In physical terms, the wave leaves its background like \(e^{mz}\), so
\(m\) is a spatial eigenvalue of the linearization about the
background, that is, a root of \(P\); and \(m\) is one of
\(1,\ldots,p\).

\paragraph{Other equations with the same profile equation.}
The reduction of Section~\ref{sec:setting} uses only the travelling-wave
variable. The counts therefore apply to any equation
\(\mathcal L u+\mathcal M(u^2)=0\) with constant-coefficient operators,
mixed derivatives and several space variables included, whose
plane-wave reduction integrates to \eqref{eq:main}; the coefficients of
\(P\) then depend on the speed and the direction of the wave.
An example in three space variables is the equation of
\citet{KudryashovSinelshchikovThreeDim2012} for waves in a liquid with
gas bubbles, whose travelling-wave reduction is that of KdV--Burgers.
The repository lists such equations by profile operator, with their
reductions (directory \texttt{tools/wave\_dictionary/}).

\paragraph{Ansatz methods.}
Many exact travelling waves in the literature are produced by
substituting a finite expansion: the tanh and sech methods, the
exp-function and \((G'/G)\)-expansion methods, the simplest-equation
method with a Riccati equation \citep{Kudryashov2005Simplest}, and the
Jacobi and Weierstrass elliptic-function expansions
\citep{Kudryashov2009Errors}. In their standard form, with one phase
\(kx-\omega t\) and constant coefficients, the first group produces
functions that are rational in one exponential and the second group
elliptic functions of one argument; rational functions of \(z\) arise
as limits. With respect to its primitive period, an output with one
pole per period is, up to the symmetries of
Proposition~\ref{prop:fixed}, one of the waves \eqref{eq:fixed-wave},
one of the elliptic waves of Section~\ref{sec:elliptic-counting}, or
the rational wave; for the equations of Theorem~\ref{thm:complete}
every nonconstant output has one pole per period. The counts therefore
serve as a checklist: at a covered order, once the rate or the period
lattice is fixed, a tanh-type or elliptic ansatz with one pole per
period can return at most \(N_p\), respectively \(M_p\), normalized
solutions, independently of the expansion method.
This agrees with the critical comparisons of these methods
\citep{KudryashovLoguinova2009,Kudryashov2009Errors,KudryashovSoukharev2009}:
the ``new'' waves they returned for the KdV--Burgers and
Kuramoto--Sivashinsky equations coincide with the known ones
\citep{Kudryashov2009KdVB,KudryashovSoukharev2009}, and an ansatz whose
pole order is lower than \(p\) returns no wave at all
\citep{Kudryashov2010Kawahara}.
For the conformable derivative \(T_\alpha f(t)=t^{1-\alpha}f'(t)\),
\(t>0\), \(0<\alpha\le1\) \citep{KhalilEtAl2014}, the change
\(\tau=t^\alpha/\alpha\) gives \(T_\alpha=\dd/\dd\tau\), and similarly
for a spatial conformable derivative, so whenever the transformed
equation has the profile equation \eqref{eq:main} its waves fall under
the same counts; this does not apply to the Caputo or
Riemann--Liouville derivatives.
